\section{The Body of Glory}

Heaven is not only for souls. The Church confesses, in the profession of faith of the Fourth Lateran Council, that all the dead \latin{cum suis propriis resurgent corporibus, quae nunc gestant}: they will rise with their own bodies, the bodies they now bear.\endnote{Fourth Lateran Council (1215), constitution 1, \emph{Firmiter credimus}, Denzinger n.~801.} The body that kneels, works, suffers and is buried is the body that will be glorified. The Apostle, asked how the dead rise and with what body, answers with the seed, which must die to live, and with four contrasts:

\begin{quote}
It is sown in corruption: it shall rise in incorruption. It is sown in dishonour: it shall rise in glory. It is sown in weakness: it shall rise in power. It is sown a natural body: it shall rise a spiritual body.\endnote{1 Cor 15:42--44, Douay--Rheims (Challoner).}
\end{quote}

The pattern of that glory is Christ risen: he ``will reform the body of our lowness, made like to the body of his glory.''\endnote{Phil 3:20--21, Douay--Rheims (Challoner).} The Church already contemplates the pattern fulfilled in one of us. Pius XII defined in 1950 that the Mother of God, ``having completed the course of her earthly life, was assumed body and soul into heavenly glory''; and Paul VI, confessing the same faith, adds that she was thereby ``likened to her risen Son in anticipation of the future lot of all the just.''\endnote{Pius XII, apostolic constitution \emph{Munificentissimus Deus} (1 November 1950), n.~44; Paul VI, \emph{Credo of the People of God} (30 June 1968), n.~15. Official English texts.} What Mary is, the Church will be.

\subsection{Why the body, if the soul already sees}

The definition of 1336 raises a question the devout have always felt. If the souls of the saints already see God and are \latin{vere beatae}, truly blessed, before the resurrection, what can the body add? Thomas Aquinas answers with great tenderness toward the body. The separated soul, he says, has the whole of its good; yet it is not wholly at rest, because it is made to be the form of a body, and it ``desires to enjoy God in such a way that the enjoyment also may overflow into the body, as far as possible.'' And so:

\begin{quote}
The desire of the separated soul is entirely at rest, as regards the thing desired; since, to wit, it has that which suffices its appetite. But it is not wholly at rest, as regards the desirer, since it does not possess that good in every way that it would wish to possess it. Consequently, after the body has been resumed, Happiness increases not in intensity, but in extent.\endnote{Aquinas, \emph{Summa theologiae} I-II, q.~4, a.~5, replies to obj.~4 and 5, trans.\ Fathers of the English Dominican Province (New York: Benziger Brothers; Project Gutenberg eBook 17897).}
\end{quote}

The vision does not get better. The person does. A human being is not a soul for a time inconvenienced by matter; a soul without its body has everything it wants and is not yet everything it is. The resurrection adds nothing to God. It completes the creature who is looking at him.

Bernard of Clairvaux, in the last chapter of his treatise on the love of God, says the same in the voice of love, in one sentence that every mourner may keep: ``when the soul leaves the body, it is gain, and greater gain when it rejoins it.'' And he adds that the soul may ``think to obtain increase of glory from its body, because, though infirm and mortal, it has contributed greatly to its merits.''\endnote{Bernard, \emph{On the Love of God} XI, trans.\ Patmore (1884).} He gives the body its honour:

\begin{quote}
The body then is for the soul a faithful companion; if it be a burden, it is also a help; when it ceases to aid, it also ceases to hinder; when again it helps, it is no more a burden. The first state is laborious, but profitable; the second inactive, but in no sense tedious; and the third is glorious.\endnote{Bernard, \emph{On the Love of God} XI, trans.\ Patmore (1884).}
\end{quote}

The soul that loves God, he says, profits by its poor, weak body ``whether living, dead, or risen---living, by bringing forth with her fruits of penance; dead, by resting from its labour; risen, the two together attain the consummation of all joy.''\endnote{Bernard, \emph{On the Love of God} XI, trans.\ Patmore (1884).} The fasts of Lent, the kneeling, the weariness of the sick bed, the work of the hands: none of it is lost. The body that served will share the reward.

\subsection{The endowments of glory}

The Doctors of the schools gathered what Scripture says of the risen body under four gifts. The Supplement to the \emph{Summa} treats them in order: impassibility, subtlety, agility and clarity. Their names say what they promise: a body that can no longer suffer harm, that is wholly at the service of the spirit, that goes where the spirit wills, and that shines.\endnote{\emph{Summa theologiae}, Supplement, qq.~82--85 (impassibility, subtlety, agility, clarity), trans.\ Fathers of the English Dominican Province, as presented at New Advent.} This is the common teaching of the Doctors, and its root is a single principle. In the mature \emph{Summa}, Thomas says that ``these gifts of a glorified body result from an overflow of the soul's glory on to the body.''\endnote{Aquinas, \emph{Summa theologiae} III, q.~28, a.~2, reply to obj.~3, trans.\ English Dominican Province (Project Gutenberg eBook 19950).} The Supplement states the principle as an order of subjection: ``the human body and all that it contains will be perfectly subject to the rational soul, even as the soul will be perfectly subject to God.''\endnote{\emph{Summa theologiae}, Supplement, q.~82, a.~1, trans.\ English Dominican Province, as presented at New Advent.}

The body's glory, then, is not a cosmetic laid on from outside. It is the vision of God shining through the flesh, as Christ shone on the mountain of the Transfiguration, and as he promised that ``the just shine as the sun, in the kingdom of their Father.''\endnote{Matt 13:43, Douay--Rheims (Challoner).} The Supplement borrows an image from Gregory the Great for that brightness: ``in the glorified body the glory of the soul will be known, even as through a crystal is known the color of a body contained in a crystal vessel.''\endnote{\emph{Summa theologiae}, Supplement, q.~85, a.~1, trans.\ English Dominican Province, as presented at New Advent.} And because the soul's glory is measured by charity, so is the body's: star differs from star here too.

Augustine, who wrote at length on the risen body, stops in the middle of his praise to confess the limit of what he knows: ``What power of movement such bodies shall possess, I have not the audacity rashly to define, as I have not the ability to conceive.'' Yet he is sure of the principle: ``the body shall forthwith be wherever the spirit wills, and the spirit shall will nothing which is unbecoming either to the spirit or to the body.'' And he is sure what the body will be for: ``All the members and organs of the incorruptible body, which now we see to be suited to various necessary uses, shall contribute to the praises of God.''\endnote{Augustine, \emph{The City of God} XXII.30, trans.\ Dods (1871).}

That is the end of the body's story. The hands that now labour will be lifted in praise; the eyes that now weep will see, in every creature of the new heavens and the new earth, the God who made them; the tongue will sing. ``This corruptible must put on incorruption: and this mortal must put on immortality. \ldots\ O death, where is thy victory? O death, where is thy sting?''\endnote{1 Cor 15:53--55, Douay--Rheims (Challoner).} The Christian keeps his body in chastity and offers it in penance because it is to be glorified; and he buries his dead with honour, as the Church does, because they will rise.
